%======================================================================
\subsection{SIM-07 P1 --- Built, 1 Nov: the Model Builder's model layer}\label{sec:sim07-model-built}
%======================================================================

The headless half of phase P1 is in \file{src/Mod/BruceWorx/bruceworx\_sim\_model.py}: no Qt, tested in FreeCADCmd.
The tree widget, the settings form and the status marks are built on it by the shell work; this layer does not touch
the solvers or the workspace files.

\paragraph{What you can do now (through the API).}
\begin{itemize}[nosep]
  \item \textbf{The tree.} \code{model\_tree(analysis)} returns Global Definitions (Parameters, Selections: empty
        placeholders), Geometry (the analysed part), Materials, Physics, Mesh, Study (a step node named for the study
        type, and the solver) and Results, in solving order. Physics lists the supports and loads of a structural
        study, the thermal conditions of a thermal one, and for an airflow study virtual nodes bound to the study's
        properties: inlet and far field, walls, turbulence and, in prop wash, one node per marked motor. Results
        lists the result objects and virtual plot nodes by study type (stress, displacement, safety factor, hoop and
        radial stress for a spinning part, modes, temperature, forces, pressure), a probe and the report. Keys are
        stable strings (\code{phys/Force}); icons are the ids of the mock-up sprite. Building the tree creates and
        changes nothing, so a study made before today opens unchanged (tested: same objects, same properties).
  \item \textbf{One settings declaration per node type.} \code{settings(node)} returns sections of fields (name,
        label, kind, unit, choices, help, range). \code{get\_value} and \code{set\_value} read and write the existing
        properties; a write is one document transaction, so one undo step. Numbers accept a quantity with a unit
        (``500\,N'', ``2\,kN'', ``1\,cm'', ``12000\,rpm'', ``20\,degC'', ``300\,K''); a unit of another kind, a value
        out of range, an unknown choice or a missing face raises a \code{ValueError} that says what was expected.
        Covered: the study (label, knockdown, stress measure), the material (a card from the library the New Study
        panel uses), Fixed, Displacement, Hold, Force (faces, magnitude, direction, reversed), Pressure, Gravity, Spin
        (speed, axis), Temperature, Heat flux, Convection, Heat source, Initial temperature, the mesh (maximum size,
        fineness, second order, node limit), the solver, the study step (modal, buckling and thermal settings only for
        their types) and the airflow groups of the old panel (speed and angle, air, far field and domain, turbulence,
        reference area and length, mesh preset, actuator disk).
  \item \textbf{Out-of-date marks.} \code{state(analysis)} gives each node ``ok'', ``stale'', ``warn'', ``todo'' or
        ``run''. \code{node\_messages(analysis)} gives the pre-check findings by node for the Messages tab.
\end{itemize}

\paragraph{Fingerprints, as in the plan.} $h(n)=H(s(n),h(\mathrm{deps}))$ with $H$ a 12-character SHA-1 of the
settings' canonical text. The geometry term is the rounded volume (6 digits), area, bounding box (5 digits), face and
solid counts of the analysed part, not \code{Shape.hashCode}, which changes at every save. The dependencies are
\[
  h_\mathrm{mesh}=H(\text{mesh settings},h_\mathrm{geo}),\qquad
  h_\mathrm{sol}=H(\text{study},\text{solver},h_\mathrm{mesh},h_\mathrm{materials},h_\mathrm{physics}),\qquad
  h_\mathrm{res}=h_\mathrm{sol}.
\]
The study keeps the fingerprints of its last mesh and last solution in a hidden \code{BW\_Fingerprints} map, added on
the first \code{record}. \code{sim.mesh\_study} records ``mesh'', \code{sim.finish\_solve} and \code{cfd.save\_summary}
record ``solution'' (one line each). A compute node is stale when its stored fingerprint differs from the current one.
Acceptance on a plate with a hole (100\,$\times$\,50\,$\times$\,5\,mm, fixed left, 500\,N right): moving the force to
other faces leaves geometry, material, supports and mesh ``ok'' and marks the solution and every result node stale;
changing the material does the same; changing the maximum element size marks the mesh and everything after it; a longer
plate marks the mesh. Save, close and reopen gives the same fingerprints and no stale mark. A real CalculiX solve
records, and a force edit afterwards marks it stale.

\paragraph{Decided while building.}
\begin{itemize}[nosep]
  \item \textbf{No local mesh sizes yet}, so in P1 the mesh depends on the geometry and its own settings only; a load's
        faces mark the mesh the day a local size refers to them (P2).
  \item \textbf{Display settings are not inputs}: the stress measure, band colours and edges, the probe point, the
        pressure colour range and the played mode scale recolour stored fields and never mark the solution. The
        knockdown is an input, because the safety factor is computed at the solve.
  \item \textbf{A study from before today} (nothing stored) counts as up to date when it has a mesh or results; the
        first re-mesh of such a study marks its old results stale, since they belong to the old mesh.
  \item \textbf{Airflow has one stored stage.} OpenFOAM meshes inside the run, so its Mesh node follows the solution.
  \item \textbf{A write never recomputes the mesh.} \code{set\_value} recomputes only the edited constraint (its
        symbols); changing the maximum size does not re-mesh, the node turns stale. Nothing solves until Compute (W6).
  \item \textbf{Findings go to the node they name} (``Force: no faces\ldots'' to the Force), the rest by keyword to the
        section (no support: Physics; no material: Materials), and what is left to the study node.
  \item \textbf{Errors mark ``todo'', warnings ``warn''}; Parameters and Selections carry no mark in P1.
\end{itemize}

\paragraph{Tests.} \file{TestBruceWorxSimModel.py}, 27 headless tests: the section order and keys (static, modal,
thermal, external flow, prop wash) and that building the tree, the state and the messages change no object and add no
property; each field type, units and their errors, one undo step, face selections and the bad-face error, material
card, spin in rpm, temperatures in K and degC, convection, the per-type step fields and the airflow fields; the
acceptance cases above, the legacy study, save and load, the running mark, the todo and warn marks and messages, and a
real CalculiX solve (skipped without \code{ccx}). \file{TestBruceWorxSimCore} and \file{TestBruceWorxSimRotor} still pass.

\paragraph{Open questions for you.}
\begin{enumerate}[nosep]
  \item \textbf{The analysed part in the Geometry node} is read only (the mesh shape object is built from it at New
        Study). Should the Settings pane let you swap it, which re-points the study and marks everything stale?
  \item \textbf{Material cards} are chosen by name from the library. Two libraries sharing a name show the library in
        brackets; is that enough, or should the pane group by library?
  \item \textbf{Hold constraints} (the spin-hold points) are recognised by their label ``Hold\ldots''. A proper type
        property on them would be safer; worth a small change in \code{hold\_for\_spin}?
\end{enumerate}
